\section{Introduction}
Anonymous credentials enable users to prove possession of attributes without revealing them, while resisting profiling and linkability. Anonymous credentials are currently gaining traction for anonymous rate limiting~\cite{yun-privacypass-crypto-arc-00,schlesinger-cfrg-act-00}, digital wallets~\cite{TODO} and a number of other applications. A key advantage of anonymous credentials is the constant issuance size. 
Achieving anonymous credentials and the underlying blind signatures in the post-quantum setting with practical sizes remains challenging. We combine a proof-friendly, lattice-based Blind Signature (vSIS-BBS~\cite{PKC:DKLW25} combined with an NTRU sampler~\cite{AC:ENSTW23}) with the SmallWood hash-based NIZK proof system~\cite{EPRINT:FenRiv25} to obtain compact, verifiable presentations. \lena{We then show how to extend SPRUCE to get the first post-quantum version of an ARC-style anonymous credential token. }

\paragraph{Contributions.}
In more detail, this work combines the following primitives:
\begin{enumerate}
    \item We implement an instantiation of the vSIS-BBS blind signature ~\cite{PKC:DKLW25}. This signature is tailored for proof systems: the message generated by the holder does not involve computing a full-fledged cryptographic hash function, simplifying the proof size towards the issuer and the verifier.
\item We adopt the trapdoor-based Gaussian preimage sampler tuned for stable per-slot tails from ~\cite{AC:ENSTW23}. This sampler generates the preimages necessary for the aforementioned signature scheme.
\item We implement the SmallWood~\cite{EPRINT:FenRiv25} NIZK proof system which is linear in the witness size, but succinct in the size of the statement being proven. We then encode the degree-2 signature relation and range membership via SmallWood tailored constraints.
\item We present implementation-driven optimizations (Merkle multiproofs, trimming, compact headers) that reduce the proof size without altering algebraic checks.
\item We combine all building blocks, together with carefully chosen parameters. This yields a Blind Signature scheme with $\approx$2\,kB signatures and $\approx$15--20\,kB proofs.
\item \lena{ARC tokens are unlinkable tokens with constant issuance size, regardless of many presentations they have. We show how do add SPRUCE to get a post-quantum ARC token in under XX kB of proof. }
\end{enumerate}

\subsection{How the pieces connect in our construction}\label{subsec:how-it-connects}

At a high level (see Fig.~\ref{fig:ac-flow}):
\begin{enumerate}[leftmargin=1.25em]
  \item \textbf{Commitment from vSIS (}\(\Holder\to\Issuer\)\textbf{).} The holder computes \(C=\Commit(\PP,m;r)\) using a vSIS-hash–based commitment (computational binding from vSIS/SIS-type assumptions; computational hiding from randomized hashing). \emph{No} attribute is revealed to \(\Issuer\).
  \item \textbf{Blind signature on \(C\) (}\(\Issuer\to\Holder\)\textbf{).} Via \(\Blind\)/\(\Sign\)/\(\Unblind\), \(\Issuer\) signs the blinded commitment; the holder obtains a signature \(\sigma\) on \(C\) without disclosing \(m\).
  \item \textbf{Zero-knowledge presentation (}\(\Holder\to\Verifier\)\textbf{).} Using SmallWood’s hash-based PCS/PIOP stack, \(\Holder\) proves knowledge of \((m,r,\sigma')\) such that \(C=\Commit(\PP,m;r)\) and \(\Verify(\mathsf{ipk},C,\sigma')=1\), and optionally that \(m\) satisfies predicates (e.g., ranges). \(\Verifier\) outputs accept/reject.
\end{enumerate}

\begin{figure}[H]
\centering
\begin{tikzpicture}[
    actor/.style={font=\bfseries, align=center},
    lifeline/.style={draw, thick},
    msg/.style={-Latex, thick},
    note/.style={draw, rounded corners, fill=gray!10, align=center, inner sep=4pt},
    tag/.style={font=\footnotesize\itshape},
    x=1.1cm, y=0.9cm
]
% --- Actors (Holder moved to the middle) ---
\node[actor] (I) at (0,0)  {\Large Issuer\\[-1pt]$\Issuer$};
\node[actor] (H) at (6,0)  {\Large Holder\\[-1pt]$\Holder$};
\node[actor] (V) at (12,0) {\Large Verifier\\[-1pt]$\Verifier$};

% --- Lifelines ---
\draw[lifeline] (I) -- ++(0,-8.2);
\draw[lifeline] (H) -- ++(0,-8.2);
\draw[lifeline] (V) -- ++(0,-8.2);

% --- Notes and messages ---
\node[note] at ($(H)+(-2.5,-1.0)$) {Blind commitment to attributes $m$ with randomness $r$};

% Commitment (Holder -> Issuer)
\draw[msg] ($(H)+(0,-2.2)$) -- node[above,sloped] {$C=\Commit(\PP,m;r)$ (blinded)} ($(I)+(0,-2.2)$);

% Signature (Issuer -> Holder)
\draw[msg] ($(I)+(0,-3.8)$) -- node[above,sloped] {Signature $\sigma$ on $C$} ($(H)+(0,-3.8)$);

\node[note] at ($(H)+(0,-4.9)$) {Credential $=(C,\sigma)$};

% Presentation (Holder -> Verifier)
\draw[msg] ($(H)+(0,-6.3)$) -- node[above,sloped] {\emph{NIZK AoK} for $(C, \sigma, m, r)$} ($(V)+(0,-6.3)$);

% Verification (Verifier -> Holder)
\draw[msg] ($(V)+(0,-7.5)$) -- node[above,sloped] {Verifier: $\Verify(\textsf{pp};\pi)$} ($(H)+(0,-7.5)$);

\end{tikzpicture}
\caption{Issuance and presentation flow. \emph{Default showing transmits only the proof~$\pi$; neither $r$ nor a deterministic tag is revealed, yielding unlinkability.}}
\label{fig:ac-flow}
\end{figure}

The public parameters for our construction are a trapdoored matrix \(\mathbf{A}\in R_q^{\,n\times m}\) (verification map) for an ISIS instance and a public matrix \(\mathbf{B}\in R_q^{\,n\times \ell}\) (hash domain).
The message and randomness are an ISIS preimage relation with a BBS-style \emph{rational hash}; its output \(\mathbf{t}\) serves as the commitment.
 For a message \(u\) and signing randomness
\(\chi\), the vSIS construction defines a rational function \(h_{u,\chi}\) and sets the target
\[
  \mathbf{t}\ \coloneqq\ h_{u,\chi}(\mathbf{B})\ \in R_q^{\,n}.
\]
A signature is precisely a short preimage \(\mathbf{s}\) such that
\[
  \boxed{\ \mathbf{A}\mathbf{s}\ \equiv\ \mathbf{t}\ \pmod q\ } \qquad\text{with}\quad \|\mathbf{s}\|\ \text{small}.
\]
This is the verification equation we will prove in zero knowledge later on.

Concrete choices of \(h_{u,\chi}\) translate the verification equations into low-degree
polynomial relations over \(R_q\)
% ; in our work, we use functions of the form
% $$h_{u,\chi}(\tilde{\mathbf{b}}_0,\tilde{\mathbf{b}}_1)
%   = \big((1,u^\top,x_0^\top)\cdot \tilde{\mathbf{b}}_0\big)\big/(\tilde{\mathbf{b}}_1-x_1)\text{,}$$
% with \(\chi=(x_0,x_1)\), which exactly yields an ISIS target \(\mathbf{t}\) when evaluated at the public
% parameter \(\mathbf{B}=(\mathbf{b}_0,\mathbf{b}_1)\).  
Previous work \cite{cryptoeprint:2025/356} formalizes this pattern through a
general signature skeleton \(\Sigma_{\mathrm{ISIS_f}}\) where verification checks \(\mathbf{A}\mathbf{s}=h_{u,\chi}(\mathbf{B})\pmod q\)
and \(\mathbf{s},\chi,u\) is short; our instantiation follows that template. 

\begin{comment}
\paragraph{Threat model (brief).}
We consider malicious holders (one-more unforgeability), malicious issuers (blindness), and verifiers (unlinkability). Security follows from SIS/ISIS hardness for signatures \cite{cryptoeprint:2025/356} and soundness/HVZK for SmallWood in ROM/FS \cite{cryptoeprint:2025/1085}.

\paragraph{Roadmap.}
Section~1 defines primitives (commitments, blind signatures, NIZK, SIS/ISIS, rational hash, vanishing polynomials). Section~2 covers algebraic preliminaries. Section~3 presents the vSIS-BBS scheme and the cleared identity. Section~4 details the Gaussian sampler. Section~5 reviews the SmallWood system (generic). Section~6 lists optimizations. Section~7 states parameter constraints and choices. Section~8 encodes vSIS-BBS in SmallWood (witness layout and constraints). Section~9 reports results.

\end{comment}